\subsection{Borel, Lebesgue, and Stieltjes Measure Constructions}

\begin{proposition} \label{prop:half-open-real-line-length-premeasure} \hlt{Half-Open Interval Algebra on $\R$}\\
Let $\mathcal{F}$ be the algebra of finite disjoint unions of $(a,b]$, $(-\infty,b]$, $(a,+\infty)$ ($a,b\in\R$), and $\R$. Set $\ell((a,b])=b-a$ and $\ell((-\infty,b])=\ell((a,+\infty))=\ell(\R)=+\infty$, extending additively over disjoint unions. Then
\begin{enumerate}[label=(\roman*)]
\setlength{\itemsep}{0pt}
\item $\mathcal{F}$ is an algebra on $\R$ and $\sigma(\mathcal{F})=\mathcal{B}(\R)$.
\item $\ell$ is a premeasure on $\mathcal{F}$.
\end{enumerate}
\begin{proof}
\begin{enumerate}[label=(\roman*)]
\setlength{\itemsep}{0pt}
\item Generator complements lie in $\mathcal{F}$, hence $\mathcal{F}$ is an algebra.\\
Since $\mathcal{F}\subseteq\mathcal{B}(\R)$ and $(a,b)=\bigcup\limits_{n>1/(b-a)}(a,b-1/n]$, $\sigma(\mathcal{F})=\mathcal{B}(\R)$.
\item Normal forms make $\ell$ well defined and finitely additive.\\
For countable additivity, reduce to $(a,b]=\bigsqcup\limits_{j\geq1}(a_j,b_j]$: finite subunions give $\sum\limits_{j=1}^{N}(b_j-a_j)\leq b-a$.\\
Conversely, choose $\varepsilon_j>0$ with $\sum\limits_{j\geq1}\varepsilon_j<\varepsilon$; by compactness, $\exists$ finite $K\subseteq\N$ s.t. $[a+\delta,b]\subseteq\bigcup\limits_{j\in K}(a_j,b_j+\varepsilon_j)$, hence $b-a-\delta\leq\sum\limits_{j\in K}(b_j-a_j+\varepsilon_j)\leq\sum\limits_{j\geq1}(b_j-a_j)+\varepsilon$.\\
Let $\delta,\varepsilon\downarrow0$. Truncation handles half-lines; finite disjoint unions follow componentwise.
\end{enumerate}
\hlt{Structural Commentary}: Lebesgue measure starts as length on a finite interval algebra, not on all subsets.\\
\hlt{Key Insight}: Compactness supplies the reverse inequality in countable additivity of interval length.
\end{proof}
\end{proposition}

\begin{remark} \label{rem:half-open-endpoints-make-an-algebra} \hlt{Why the Half-Open Convention Matters}\\
Finite unions generated by $(a,b)$, $(-\infty,b)$, $(a,+\infty)$, and $\R$ still generate $\mathcal{B}(\R)$, but are not an algebra, where\\
$(0,1)^c=(-\infty,0]\cup[1,+\infty)$. Half-open endpoints make complements remain in the starting family.
\end{remark}

\begin{definition} \label{def:borel-lebesgue-stieltjes-measures} \hlt{Borel--Stieltjes and Lebesgue--Stieltjes Measures}\\
Let $g:\R\to\R$ be increasing and right-continuous, and set $g(\pm\infty)\defeq\lim_{x\to\pm\infty}g(x)\in\overline{\R}$.\\
Let $\mathcal{I}\defeq\{(a,b]:a<b\}\cup\{(-\infty,b]:b\in\R\}\cup\{(a,+\infty):a\in\R\}\cup\{\R\}$ and $\mathcal{F}\defeq\operatorname{alg}(\mathcal{I})$.
\begin{enumerate}[label=(\roman*)]
\setlength{\itemsep}{0pt}
\item Define $G_g$ on $\mathcal{I}$ by endpoint increments and set $\mu_g(\bigsqcup_{j=1}^{m}I_j)\defeq\sum_{j=1}^{m}G_g(I_j)$ for disjoint $I_j\in\mathcal{I}$.
\item The unique extension of $\mu_g$ to $\mathcal{B}(\R)$ is the \hlt{Borel--Stieltjes measure generated by $g$}.
\item The completion of this Borel measure is the \hlt{Lebesgue--Stieltjes measure generated by $g$}.
\end{enumerate}
\end{definition}

\begin{theorem} \label{thm:borel-stieltjes-measure-continuous-increasing}\label{thm:borel-stieltjes-measure-generation} \hlt{Borel--Stieltjes Measure Existence}\\
Let $g:\R\to\R$ be increasing and right-continuous, and put $g(\pm\infty)\defeq\lim_{x\to\pm\infty}g(x)\in\overline{\R}$.\\
Set $\mathcal{I}\defeq\{(a,b]:a<b\}\cup\{(-\infty,b]:b\in\R\}\cup\{(a,+\infty):a\in\R\}\cup\{\R\}$ and $\mathcal{F}\defeq\operatorname{alg}(\mathcal{I})$.\\
Define $G_g$ on $\mathcal{I}$ by endpoint increments and set $\mu_g(\bigsqcup_{j=1}^{m}I_j)\defeq\sum_{j=1}^{m}G_g(I_j)$ for disjoint $I_j\in\mathcal{I}$. Then:
\begin{enumerate}[label=(\roman*)]
\setlength{\itemsep}{0pt}
\item $\mu_g$ is a $\sigma$-finite premeasure on $\mathcal{F}$ and has a unique extension to $\mathcal{B}(\R)$, again denoted $\mu_g$.
\item If $g=\operatorname{id}_{\R}$, then $\mu_g$ is Lebesgue measure.
\item If $g$ is continuous, then $\mu_g(J)=g(b)-g(a)$ for $J\in\{(a,b),[a,b),(a,b],[a,b]\}$, $a<b$.
\end{enumerate}
\begin{proof}
\begin{enumerate}[label=(\roman*)]
\setlength{\itemsep}{0pt}
\item Normal forms give well-defined finite additivity: $G_g$-increments telescope and $G_g\geq0$.\\
If $(a,b]=\bigsqcup_{j\geq1}(a_j,b_j]$, finite subunions give $\sum_{j=1}^{N}G_g((a_j,b_j])\leq G_g((a,b])$.\\
Conversely, compactness gives finite $K$ with $[a+\varepsilon,b]\subseteq\bigcup_{j\in K}(a_j,b_j+\varepsilon_j)$. Hence\\
$g(b)-g(a+\varepsilon)\leq\sum_{j\in K}(g(b_j+\varepsilon_j)-g(a_j))$. Let $\varepsilon,\varepsilon_j\downarrow0$, hence countable additivity on bounded generators.\\
Half-lines follow by truncation, finite unions componentwise, so $\mu_g$ is a premeasure. Since $\R=\bigcup_{n\geq1}(-n,n]$ and $\mu_g((-n,n])=g(n)-g(-n)<+\infty$, $\mu_g$ is $\sigma$-finite; Theorem \ref{thm:sigma-finite-generated-extension-uniqueness} gives uniqueness.
\item If $g=\operatorname{id}_{\R}$, then $G_g((a,b])=b-a$, so $\mu_g$ agrees with half-open length on $\mathcal{F}$.\\
Uniqueness identifies its Borel extension with Lebesgue measure.
\item If $g$ is continuous, then $(x-1/n,x]\downarrow\{x\}$ gives $\mu_g(\{x\})=\lim_n(g(x)-g(x-1/n))=0$.
\end{enumerate}
\hlt{Structural Commentary}: Stieltjes measure replaces ordinary length $b-a$ by the increment $g(b)-g(a)$.\\
\hlt{Key Insight}: Right-continuity aligns the half-open interval convention with countable covers.
\end{proof}
\end{theorem}

\begin{proposition} \label{prop:finite-borel-measure-distribution-function} \hlt{Finite Borel Measures Give Stieltjes Generators}\\
Let $\nu$ be finite on $\mathcal{B}(\R)$ and set $g(x)\defeq\nu((-\infty,x])$. Then:
\begin{enumerate}[label=(\roman*)]
\setlength{\itemsep}{0pt}
\item $g$ is increasing and right-continuous;
\item $g(-\infty)=0$ and $g(+\infty)=\nu(\R)$;
\item $\nu((a,b])=g(b)-g(a)$ for $a\leq b$;
\item $\nu=\mu_g$ on $\mathcal{B}(\R)$.
\end{enumerate}
\begin{proof}
\begin{enumerate}[label=(\roman*)]
\setlength{\itemsep}{0pt}
\item If $a\leq b$, then $(-\infty,a]\subseteq(-\infty,b]$, so $g(a)\leq g(b)$. If $x_n\downarrow x$, then $(-\infty,x_n]\downarrow(-\infty,x]$.\\
Finite continuity from above gives $g(x_n)\downarrow g(x)$.
\item Since $(-\infty,-n]\downarrow\emptyset$ and $(-\infty,n]\uparrow\R$, continuity from above/below gives $g(-\infty)=0$ and $g(+\infty)=\nu(\R)$.
\item For $a\leq b$, $(-\infty,b]=(-\infty,a]\sqcup(a,b]$, hence $\nu((a,b])=g(b)-g(a)$.
\item By (ii)--(iii), $\nu$ and $\mu_g$ agree on $\mathcal{I}$, hence on $\mathcal{F}$. Theorem \ref{thm:borel-stieltjes-measure-generation} gives $\nu=\mu_g$ on $\mathcal{B}(\R)$.
\end{enumerate}
\hlt{Structural Commentary}: A finite real-line measure is encoded by its cumulative mass function.\\
\hlt{Key Insight}: Right-continuity of distribution functions is continuity from above for measures.
\end{proof}
\end{proposition}

\begin{example} \label{ex:borel-stieltjes-basic-generators} \hlt{Basic Borel--Stieltjes Generators}\\
The following generators illustrate the interval rule in Theorem \ref{thm:borel-stieltjes-measure-generation}.
\begin{enumerate}[label=(\roman*)]
\setlength{\itemsep}{0pt}
\item If $g(x)=2x$, then $\mu_g=2m$ on $\mathcal{B}(\R)$.
\item If $g(x)=\arctan x$, then $\mu_g((a,b])=\arctan b-\arctan a=\int_a^b(1+x^2)^{-1}\,dx$ and $\mu_g(\R)=\pi$.
\item If $g(x)=0$ for $x<0$ and $g(x)=1$ for $x\geq0$, then $\mu_g=\delta_0$.
\item If $g(x)=0$ for $x<0$ and $g(x)=x$ for $x\geq0$, then $\mu_g(E)=m(E\cap[0,+\infty))$ for Borel $E$.
\end{enumerate}
\end{example}

\begin{lemma} \label{lem:open-sets-countable-unions-rational-boxes} \hlt{Open Sets Are Countable Unions of Rational Boxes}\\
For every open $G\subseteq\R^n$, there is a countable $\Sigma\subseteq\{\prod_{i=1}^{n}(a_i,b_i):a_i,b_i\in\Q,\ a_i<b_i\}$ such that $G=\bigcup_{B\in\Sigma}B$.
\begin{proof}
Let $\mathcal{B}_{\Q}\defeq\{\prod_{i=1}^{n}(a_i,b_i):a_i,b_i\in\Q,\ a_i<b_i\}$; then $\mathcal{B}_{\Q}$ is countable since $\Q^{2n}$ is countable.\\
Set $\Sigma\defeq\{B\in\mathcal{B}_{\Q}:B\subseteq G\}$. The inclusion $\bigcup_{B\in\Sigma}B\subseteq G$ is immediate.\\
If $x\in G$, choose $r>0$ with $B_r(x)\subseteq G$, then choose $a_i,b_i\in\Q$ with $x_i-r/n<a_i<x_i<b_i<x_i+r/n$.\\
Thus $x\in B\defeq\prod_i(a_i,b_i)\subseteq B_r(x)\subseteq G$, so $x\in\bigcup_{B\in\Sigma}B$.\\
\hlt{Structural Commentary}: Rational boxes give a countable Euclidean base.\\
\hlt{Key Insight}: Openness supplies a ball; density of $\Q$ supplies rational box endpoints inside it.
\end{proof}
\end{lemma}

\begin{theorem} \label{thm:lebesgue-measure-basic-properties-checklist} \hlt{Lebesgue Measure Basic Properties}\\
Let $m=m_n^\ast|_{\mathcal{L}^n}$.
\begin{enumerate}[label=(\roman*)]
\setlength{\itemsep}{0pt}
\item $\mathcal{B}(\R^n)\subseteq\mathcal{L}^n$, and $\mathcal{L}^n$ is closed under $E\mapsto E^c$, countable $\cup$, and countable $\cap$.
\item For $E,F,E_k\in\mathcal{L}^n$, $0\leq m(E)$, $E\subseteq F\Rightarrow m(E)\leq m(F)$, and $m(\bigcup_kE_k)\leq\sum_km(E_k)$.
\item If $(E_k)\subseteq\mathcal{L}^n$ are pairwise disjoint, then $m(\bigsqcup_kE_k)=\sum_km(E_k)$.
\item $m([0,1]^n)=1$, and $E\in\mathcal{L}^n$, $t\in\R^n$ imply $E+t\in\mathcal{L}^n$ and $m(E+t)=m(E)$.
\end{enumerate}
\begin{proof}
\begin{enumerate}[label=(\roman*)]
\setlength{\itemsep}{0pt}
\item By Theorem \ref{thm:caratheodory-measurable-sets-sigma-algebra}, $\mathcal{L}^n$ is a $\sigma$-algebra, hence closed under complements and countable unions/intersections. Open boxes are cells, hence in $\mathcal{L}^n$ by Proposition \ref{prop:cells-lebesgue-measurable-volume}; by Lemma \ref{lem:open-sets-countable-unions-rational-boxes}, each open $G=\bigcup_jB_j$ with rational open boxes $B_j$, so $G\in\mathcal{L}^n$. Thus $\mathcal{B}(\R^n)=\sigma(\mathcal{O}_{\R^n})\subseteq\mathcal{L}^n$.
\item Since $m=m_n^\ast|_{\mathcal{L}^n}$, nonnegativity follows from $m_n^\ast\geq0$. If $E\subseteq F$, then $m(E)\leq m(F)$ by outer-measure monotonicity. Also $m(\bigcup_kE_k)=m_n^\ast(\bigcup_kE_k)\leq\sum_km_n^\ast(E_k)=\sum_km(E_k)$ by outer-measure subadditivity.
\item If $(E_k)\subseteq\mathcal{L}^n$ pairwise disjoint, by Theorem \ref{thm:caratheodory-measurable-sets-sigma-algebra}, $m$ countably additive on $\mathcal{L}^n$, so $m(\bigsqcup_kE_k)=\sum_km(E_k)$.
\item By Theorem \ref{thm:outer-measure-agrees-with-cell-volume}, since $[0,1]^n$ is a cell, $m([0,1]^n)=\operatorname{vol}_n([0,1]^n)=1$.\\
For translation, let $A\subseteq\R^n$. By Theorem \ref{thm:outer-measure-translation-invariance} and\\
$E\in\mathcal{L}^n$, $m_n^\ast(A)=m_n^\ast(A-t)=m_n^\ast((A-t)\cap E)+m_n^\ast((A-t)\cap E^c)=m_n^\ast(A\cap(E+t))+m_n^\ast(A\cap(E+t)^c)$; hence $E+t\in\mathcal{L}^n$, and $m(E+t)=m_n^\ast(E+t)=m_n^\ast(E)=m(E)$.
\end{enumerate}
\hlt{Structural Commentary}: The wish list is recovered only after restricting outer measure to measurable sets.\\
\hlt{Key Insight}: The box-cover construction gives size; the measurable-set condition restores additivity.
\end{proof}
\end{theorem}

\begin{theorem} \label{thm:lebesgue-measure-unique-on-open-cells} \hlt{Uniqueness from Open Cells}\\
Let $\nu$ be a measure on $\mathcal{L}^n$ with $\nu(I)=\operatorname{vol}_n(I)$ for every open cell $I\subseteq\R^n$. Then $\nu=m$ on $\mathcal{L}^n$.
\begin{proof}
Fix $N\in\N$, $Q_N=(-N,N)^n$. If $E\in\mathcal{L}^n$ and $E\subseteq Q_N$, then any open-cell cover $E\subseteq\bigcup_jJ_j$ gives $\nu(E)\leq\sum_j\nu(J_j)=\sum_j\operatorname{vol}_n(J_j)$, hence $\nu(E)\leq m(E)$ by infimum.\\
Applying this to $E$ and $Q_N\setminus E$, and using $\nu(Q_N)=m(Q_N)<+\infty$, yields $\nu(E)+\nu(Q_N\setminus E)=m(E)+m(Q_N\setminus E)$ with each left summand $\leq$ its right counterpart; hence $\nu(E)=m(E)$.\\
For arbitrary $E\in\mathcal{L}^n$, set $A_1=E\cap Q_1$, $A_N=E\cap(Q_N\setminus Q_{N-1})$ for $N>1$.\\
Then $E=\bigsqcup_NA_N$, each $A_N\subseteq Q_N$, and countable additivity gives $\nu(E)=\sum_N\nu(A_N)=\sum_Nm(A_N)=m(E)$.\\
\hlt{Structural Commentary}: Agreement on coordinate cells determines the whole Lebesgue measure.\\
\hlt{Key Insight}: Covers give $\nu\leq m$ locally; finite complements force equality.
\end{proof}
\end{theorem}

\begin{method} \label{meth:translation-invariance-caratheodory} \hlt{Translation Invariance by Moving the Test Set}
\begin{enumerate}[label=(\roman*)]
\setlength{\itemsep}{0pt}
\item For $t\in\R^n$ and $A\subseteq\R^n$, translated box covers give $m_n^\ast(A+t)=m_n^\ast(A)$.
\item If $E\in\mathcal{L}^n$, then for all $A\subseteq\R^n$, $(A\cap(E+t))-t=(A-t)\cap E$ and $(A\cap(E+t)^c)-t=(A-t)\cap E^c$.
\item Apply the Caratheodory split for $E$ to $A-t$ to get $E+t\in\mathcal{L}^n$ and $m(E+t)=m(E)$.
\end{enumerate}
\hlt{Key Insight}: Translation invariance of outer measure transports the clean-splitting identity.
\end{method}

\begin{proposition} \label{prop:borel-hierarchy-first-examples} \hlt{First Borel Hierarchy Examples}\\
In $\R^p$:
\begin{enumerate}[label=(\roman*)]
\setlength{\itemsep}{0pt}
\item If $G=\bigcap\limits_{n\geq1}U_n$ with $U_n$ open or $F=\bigcup\limits_{n\geq1}C_n$ with $C_n$ closed, then $G,F\in\mathcal{B}(\R^p)\subseteq\mathcal{L}^p$.
\item Any set obtained from open or closed sets by countably many $\cup,\cap,{}^c$ operations is Borel.
\item In $\R$, $\{0\}=\bigcap\limits_{n\geq1}(-\frac{1}{n},\frac{1}{n})$ is $G_\delta$ but not open.
\item In $\R$, $(0,1)=\bigcup\limits_{n\geq2}[\frac{1}{n},1-\frac{1}{n}]$ is $F_\sigma$ but not closed.
\end{enumerate}
\begin{proof}
\begin{enumerate}[label=(\roman*)]
\setlength{\itemsep}{0pt}
\item Since $\mathcal{B}(\R^p)=\sigma(\mathcal{O}_{\R^p})$, it contains open sets, closed sets by complements, and countable $\cap,\cup$ by $\sigma$-algebra closure; Theorem \ref{thm:lebesgue-measure-basic-properties-checklist} gives $\mathcal{B}(\R^p)\subseteq\mathcal{L}^p$.
\item This is Proposition \ref{prop:sigma-algebra-closure-and-generation} applied to $\mathcal{B}(\R^p)$.
\item For $x\in\R$, $x=0\iff \abs{x}<\frac{1}{n}$ for all $n\geq1$; the singleton $\{0\}$ is not open.
\item For $x\in\R$, $0<x<1\iff\exists n\geq2:\frac{1}{n}\leq x\leq1-\frac{1}{n}$; the interval $(0,1)$ is not closed.
\end{enumerate}
\hlt{Structural Commentary}: Borel sets are generated by countable topological operations.\\
\hlt{Key Insight}: Countable intersections and unions quickly leave the class of open and closed sets.
\end{proof}
\end{proposition}

\begin{perspective} \label{per:outer-measure-not-enough} \hlt{Why Not Measure Every Subset?}\\
Vitali-type sets show $m_1^\ast$ is not additive on $2^\R$: disjoint unions and refinements can break additivity.\\
Thus measure is $m^\ast$ on clean splitters $E$: $m^\ast(X)=m^\ast(X\cap E)+m^\ast(X\cap E^c)$ for all $X$.\\
\hlt{Key Insight}: Outer measure is universal; measure lives on clean splitters.
\end{perspective}

\begin{definition} \label{def:null-sets-and-real-line-measure-zero} \hlt{Null Sets and Real-Line Measure Zero}\\
Let $(X,\mathcal{A},\mu)$ be a measure space and let $A\subseteq\R$.
\begin{enumerate}[label=(\roman*)]
\setlength{\itemsep}{0pt}
\item $N\subseteq X$ is \hlt{null} iff $N\in\mathcal{A}$ and $\mu(N)=0$.
\item $A$ has \hlt{measure zero} iff $\forall\varepsilon>0$, $\exists$ open intervals $(O_n)$ with $A\subseteq\bigcup_{n=1}^{\infty}O_n$ and $\sum_{n=1}^{\infty}\abs{O_n}\leq\varepsilon$.
\item If $A\subseteq[a,b]$, then $A$ has \hlt{content zero} iff the condition in (ii) holds with finitely many intervals.
\end{enumerate}
For Lebesgue measure on $\R$, null sets are also \hlt{zero sets}.
\end{definition}

\begin{proposition} \label{prop:outer-null-sets-lebesgue-measurable} \hlt{Outer-Null Sets Are Lebesgue Measurable}\\
If $Z\subseteq\R^n$ and $m_n^\ast(Z)=0$, then:
\begin{enumerate}[label=(\roman*)]
\setlength{\itemsep}{0pt}
\item $Z\in\mathcal{L}^n$;
\item $m(Z)=0$;
\item $\mathcal{P}(Z)\subseteq\mathcal{L}^n$ and $m(W)=0$ for every $W\subseteq Z$.
\end{enumerate}
\begin{proof}
\begin{enumerate}[label=(\roman*)]
\setlength{\itemsep}{0pt}
\item For $A\subseteq\R^n$, monotonicity gives $m_n^\ast(A\cap Z)=0$, so $m_n^\ast(A)\leq m_n^\ast(A\cap Z)+m_n^\ast(A\cap Z^c)=m_n^\ast(A\cap Z^c)\leq m_n^\ast(A)$; hence $Z$ satisfies Caratheodory's criterion.
\item Since $Z\in\mathcal{L}^n$, $m(Z)=m_n^\ast(Z)=0$.
\item If $W\subseteq Z$, then $m_n^\ast(W)\leq m_n^\ast(Z)=0$; applying (i)--(ii) to $W$ gives $W\in\mathcal{L}^n$ and $m(W)=0$.
\end{enumerate}
\hlt{Structural Commentary}: Lebesgue measure is complete because outer-null subsets split every test set trivially.\\
\hlt{Key Insight}: A zero piece cannot disturb additivity.
\end{proof}
\end{proposition}

\begin{proposition} \label{prop:finite-countable-cantor-null-sets} \hlt{Basic Measure-Zero Examples on $\R$}
\begin{enumerate}[label=(\roman*)]
\setlength{\itemsep}{0pt}
\item Every finite $A\subseteq\R$ has content $0$.
\item Every countable $A\subseteq\R$ has measure $0$.
\item The Cantor set $P$ has content $0$.
\end{enumerate}
\begin{proof}
\begin{enumerate}[label=(\roman*)]
\setlength{\itemsep}{0pt}
\item For $A=\{a_1,\ldots,a_N\}$ and $\varepsilon>0$, choose open $I_j\ni a_j$ with $\abs{I_j}<\frac{\varepsilon}{N}$; then $A\subseteq\bigcup\limits_{j=1}^{N}I_j$ and $\sum\limits_{j=1}^{N}\abs{I_j}<\varepsilon$. The case $A=\emptyset$ is trivial.
\item For $A=\{a_j:j\in\N\}$ and $\varepsilon>0$, choose open $I_j\ni a_j$ with $\abs{I_j}<\frac{\varepsilon}{2^j}$; then $A\subseteq\bigcup\limits_{j\geq1}I_j$ and $\sum\limits_{j\geq1}\abs{I_j}<\varepsilon$.
\item Given $\varepsilon>0$, choose $N$ with $\left(\frac{2}{3}\right)^N<\frac{\varepsilon}{2}$. If $E_N=\bigcup_{k=1}^{2^N}J_k$, $\abs{J_k}=3^{-N}$, choose open $I_k\supseteq J_k$ with $\abs{I_k}<3^{-N}+\frac{\varepsilon}{2^{N+1}}$; then $P\subseteq\bigcup_{k=1}^{2^N}I_k$ and $\sum_{k=1}^{2^N}\abs{I_k}<\left(\frac{2}{3}\right)^N+\frac{\varepsilon}{2}<\varepsilon$.
\end{enumerate}
\hlt{Structural Commentary}: Zero size means arbitrarily small interval covers.\\
\hlt{Key Insight}: Lists split $\varepsilon$ by terms; Cantor stages have length $(\frac{2}{3})^N\to0$.
\end{proof}
\end{proposition}

\begin{corollary} \label{cor:uncountable-lebesgue-null-sets} \hlt{Uncountable Null Sets Exist}\\
There exists $P\subseteq\R$ such that:
\begin{enumerate}[label=(\roman*)]
\setlength{\itemsep}{0pt}
\item $\abs{P}=2^{\aleph_0}$;
\item $P\in\mathcal{L}^1$ and $m(P)=0$.
\end{enumerate}
\begin{proof}
Let $P\defeq\{\sum_{k=1}^{\infty}\varepsilon_k3^{-k}:\varepsilon_k\in\{0,2\}\}$.
\begin{enumerate}[label=(\roman*)]
\setlength{\itemsep}{0pt}
\item The map $(\varepsilon_k)\mapsto\sum_k\varepsilon_k3^{-k}$ identifies $P$ with $\{0,2\}^{\N}$, so $\abs{P}=2^{\aleph_0}$.
\item At stage $N$, $P$ is covered by $2^N$ closed intervals of length $3^{-N}$, hence $m_1^\ast(P)\leq(2/3)^N$ for all $N$.\\
Letting $N\to\infty$ gives $m_1^\ast(P)=0$, so outer-null completeness gives $P\in\mathcal{L}^1$ and $m(P)=0$.
\end{enumerate}
\hlt{Structural Commentary}: Cantor coding separates cardinality from length.\\
\hlt{Key Insight}: Null need not mean countable.
\end{proof}
\end{corollary}

\begin{corollary} \label{cor:real-line-lebesgue-measure-checks} \hlt{Real-Line Lebesgue Measure Checks}\\
Let $m$ be Lebesgue measure on $\R$.
\begin{enumerate}[label=(\roman*)]
\setlength{\itemsep}{0pt}
\item $a<b\Rightarrow m((a,b))=m((a,b])=m([a,b))=m([a,b])=b-a$.
\item $A\subseteq\R$ countable $\Rightarrow A\in\mathcal{B}(\R)$ and $m(A)=0$.
\item $\emptyset\neq G\subseteq\R$ open $\Rightarrow m(G)>0$.
\item $K\subseteq\R$ compact $\Rightarrow m(K)<+\infty$.
\end{enumerate}
\begin{proof}
\begin{enumerate}[label=(\roman*)]
\setlength{\itemsep}{0pt}
\item By cell volume, $m([a,b])=b-a$; endpoint sets finite null; add or delete $\{a\},\{b\}$ does not change measure.
\item If $A=\{a_k:k\in J\}$ with $J$ finite or countable, then $A=\bigcup_{k\in J}\{a_k\}\in\mathcal{B}(\R)$; $m(A)\leq\sum_{k\in J}m(\{a_k\})=0$.
\item Choose $(a,b)\subseteq G$ with $a<b$; by monotonicity and (i), $m(G)\geq m((a,b))=b-a>0$.
\item Compact $K$ is bounded, so $K\subseteq[-R,R]$ for some $R>0$; hence $m(K)\leq m([-R,R])=2R<+\infty$.
\end{enumerate}
\hlt{Structural Commentary}: On $\R$, length, null endpoints, and compact boundedness give quick measure estimates.\\
\hlt{Key Insight}: Nonempty open sets contain intervals; compact sets fit inside one.
\end{proof}
\end{corollary}

\begin{example} \label{ex:fat-cantor-set-positive-measure-empty-interior} \hlt{A Positive-Measure Cantor-Type Set with No Interval}\\
Let $S_0=[0,1]$ and, for $n\geq1$, let $R_n$ be the union of the middle open intervals of length $2^{-2n}$ removed from the $2^{n-1}$ components of $S_{n-1}$; set $S_n\defeq S_{n-1}\setminus R_n$ and $S\defeq\bigcap_{n=0}^{\infty}S_n$.\\
Then $m(R_n)=2^{-n-1}$, $\sum_{n\geq1}m(R_n)=1/2$, so $S=[0,1]\setminus\bigcup_nR_n$ is compact with $m(S)=1/2$.\\
Since every component of $S_n$ has length $\leq2^{-n}$, no interval $(a,b)$ can lie in $S$; hence $\operatorname{int}S=\emptyset$.\\
\hlt{Structural Commentary}: Summable removals leave measure; repeated splitting kills interior.\\
\hlt{Key Insight}: Total removed length controls measure; repeated splitting controls interior.
\end{example}
